The replicated calculations establish a positive excess-entropy--mobility association in the selected solvent-design families, with clear limits on its interpretation.
\begin{enumerate}
\item \textbf{Component count alone has no demonstrated transport effect.} An exact label-only split raises formal ideal mixing entropy while leaving diffusion identical. Monovalent salt addition slightly lowers excess entropy and gives no resolved divalent speedup.
\item \textbf{HNC-guided interaction changes produce higher-excess-entropy, faster-ion liquids.} MD confirms entropy increases of approximately $0.30k_B$ and divalent diffusion gains of 18--20\% relative to the original baselines. Both signs separately increase. These are successful design comparisons, not isolated entropy interventions.
\item \textbf{Controlled diversification gives a weaker, provisional ion response.} At fixed charged composition and mean bare pair amplitude, excess entropy rises by approximately $0.103k_B$, mean diffusion by 8--9\%, and probe diffusion by approximately 5\%. The ion significance is sensitive to fitting window and unresolved with monovalent background. Longer independent sampling and the B3 charge-sign convergence check remain necessary.
\item \textbf{Rosenfeld-like trends extend to both probe signs but do not establish a universal law.} Separate fits of $D_{+2}$ and $D_{-2}$ against either TI entropy or MD pair entropy have positive slopes and high descriptive $R^2$ within the selected families. Three or four composition means cannot determine universality, and changing interactions also changes local friction and solvation. The matched U6/H7 probe contrast remains unresolved. Species-resolved kinetics, box-size comparisons, and state-point validation are the next discriminating tests.
\item \textbf{Corrected 2PT is internally consistent but not a quantitative TI replacement.} All 24 records pass analytic/reference and mechanical-label checks, while state-mean probe spectral diffusion agrees with MSD within $2.6\%$. Entropy biases nevertheless vary with composition: the controlled monovalent-background increment is overestimated by approximately $35\%$ in the primary model. The partition sensitivity does not consistently remove this error. We retain TI for the entropy--mobility comparison and use spectra as supporting dynamics, not an independent scaling validation or a literal free-ion fraction.
\item \textbf{MD pair entropy captures most, but not all, of the excess entropy.} The all-species $s_2$ accounts for approximately 80--82\% of TI $s_{\rm ex}$ across the eight states and 79--82\% of the controlled entropy increments. Substituting $s_2$ alone would preserve the qualitative entropy ordering and positive descriptive diffusion trends, while understating the increments and retaining a composition-dependent beyond-pair residual of about $-0.32$ to $-0.36k_B$ per particle. The pair approximation does not strengthen the causal or universal interpretation of the probe response.
\end{enumerate}
